\documentclass[11pt,a4paper]{article}
\usepackage[italian]{babel}
\usepackage[margin=2.5cm]{geometry}
\usepackage{parskip}
\usepackage{amsmath, amssymb, amsthm}
\usepackage{booktabs}
\usepackage{array}
\usepackage{xcolor}
\usepackage{needspace}
\usepackage{enumitem}
\usepackage{listings}
\usepackage{tikz}
\usepackage[most]{tcolorbox}
\usepackage[hidelinks]{hyperref}
\usetikzlibrary{decorations.pathreplacing, shapes.geometric, arrows.meta, positioning}

% Niente vedove/orfane: i paragrafi non lasciano righe isolate a fine/inizio pagina
\widowpenalty=10000
\clubpenalty=10000

% Elenchi più compatti
\setlist{itemsep=2pt, parsep=0pt, topsep=4pt, partopsep=0pt}

% --- Riquadri con bordo nero, senza numero (si spezzano tra due pagine) ---
% Uso: \begin{definizione} ... \end{definizione}
%      \begin{teorema}[Nome] ... \end{teorema}  (il nome è facoltativo)
\tcbset{nota/.style={enhanced jigsaw, breakable, colback=white,
  colframe=black, boxrule=0.5pt, sharp corners}}
\NewTColorBox{definizione}{}{nota,
  before upper={\textbf{Definizione.}\ }}
\NewTColorBox{teorema}{o}{nota,
  before upper={\textbf{Teorema\IfValueT{#1}{ (#1)}.}\ }}
\NewTColorBox{proposizione}{o}{nota,
  before upper={\textbf{Proposizione\IfValueT{#1}{ (#1)}.}\ }}
\NewTColorBox{proprieta}{o}{nota,
  before upper={\textbf{Proprietà\IfValueT{#1}{ (#1)}.}\ }}
\NewTColorBox{assioma}{o}{nota, boxrule=1pt,
  before upper={\textbf{Assioma\IfValueT{#1}{ (#1)}.}\ }}

% --- Ambienti semplici (senza riquadro) ---
% Il titolo sta in cima al blocco, anche se il contenuto inizia con un riquadro o
% con codice; e non resta mai da solo in fondo a una pagina.
% Uso: \begin{esempio}[Titolo facoltativo] ... \end{esempio}
\newcommand{\newrunin}[2]{%
  \NewDocumentEnvironment{#1}{o}{%
    \par\Needspace{4\baselineskip}\addvspace{\medskipamount}\noindent
    \textbf{#2}\IfValueT{##1}{ (##1)}\textbf{.}\ \ignorespaces}%
  {\par\addvspace{\medskipamount}}}
\newrunin{esempio}{Esempio}
\newrunin{esercizio}{Esercizio}
\newrunin{osservazione}{Osservazione}

% V e F nelle tavole di verità
\newcommand{\V}{\textsf{V}}
\newcommand{\F}{\textsf{F}}

% --- Codice C ---
% Uso: \begin{codice} ... \end{codice}      codice C numerato
%      \begin{comando} ... \end{comando}    comandi da terminale
%      \begin{risultato} ... \end{risultato} output di un programma
%      \lstinline|printf("ciao");|           codice dentro al testo
\lstdefinestyle{c}{
  language=C,
  basicstyle=\ttfamily\small,
  keywordstyle=\color{blue}\bfseries,
  commentstyle=\color{gray}\itshape,
  stringstyle=\color{red!70!black},
  numbers=left,
  numberstyle=\tiny\color{gray},
  numbersep=8pt,
  columns=flexible,
  keepspaces=true,
  breaklines=true,
  showstringspaces=false,
  tabsize=4,
  morekeywords={include, define},
  % lettere accentate nei commenti
  literate={à}{{\`a}}1 {è}{{\`e}}1 {é}{{\'e}}1 {ì}{{\`i}}1
           {ò}{{\`o}}1 {ù}{{\`u}}1 {È}{{\`E}}1 {À}{{\`A}}1
}
\lstset{style=c}
\newtcblisting{codice}{listing only, nota, left=6mm,
  listing options={style=c}}
\newtcblisting{comando}{listing only, nota,
  listing options={basicstyle=\ttfamily\small, numbers=none, language={},
    columns=flexible, keepspaces=true, breaklines=true}}
\newtcblisting{risultato}{listing only, enhanced jigsaw, breakable,
  colback=black!5, colframe=black, boxrule=0.5pt, sharp corners,
  listing options={basicstyle=\ttfamily\small, numbers=none, language={},
    columns=flexible, keepspaces=true, breaklines=true}}

% --- Diagrammi di flusso (simboli standard) ---
\tikzset{
  terminale/.style = {ellipse, draw, minimum width=2.5cm, minimum height=0.9cm},
  io/.style        = {trapezium, trapezium left angle=70, trapezium right angle=110,
                      draw, minimum width=2.5cm, minimum height=0.9cm},
  processo/.style  = {rectangle, draw, minimum width=2.5cm, minimum height=0.9cm},
  decisione/.style = {diamond, draw, aspect=2, minimum width=2.5cm},
  freccia/.style   = {thick, -{Stealth}}
}

\title{Radici di polinomi e ricorsione lineare}
\author{Marco Cerbella}
\date{Lezione 6 -- 6 ottobre 2026}

\begin{document}
\maketitle

\section{Radici di polinomi}

\begin{definizione}
Sia $f \in \mathbb{R}[t]$. Un numero $\alpha \in \mathbb{R}$ si dice \emph{radice}
di $f$ se $f(\alpha) = 0$.
\end{definizione}

Il legame con la divisibilità, vista nella lezione precedente, è il seguente.

\begin{teorema}[del fattore]
Sia $f \in \mathbb{R}[t]$ e sia $\alpha \in \mathbb{R}$. Allora
\[
  \alpha \text{ è una radice di } f
  \iff
  (t - \alpha) \mid f .
\]
\end{teorema}

\begin{proof}
Dividiamo $f$ per $t - \alpha$ (che ha grado $1$): per il teorema di divisione
con resto esistono $q, r \in \mathbb{R}[t]$ tali che
\[
  f(t) = q(t)\,(t - \alpha) + r(t), \qquad \operatorname{gr}(r) < 1 ,
\]
quindi $r$ è un polinomio costante. Valutando in $t = \alpha$ si ottiene
\[
  f(\alpha) = q(\alpha)\cdot 0 + r = r .
\]
Dunque $f(\alpha) = 0$ se e solo se $r = 0$, cioè se e solo se $t - \alpha$
divide $f$.
\end{proof}

\begin{osservazione}
Il resto della divisione di $f$ per $t - \alpha$ è quindi esattamente il valore
$f(\alpha)$. Se $\alpha$ è radice di $f$ si può scrivere $f = (t - \alpha)\, q$,
con $\operatorname{gr}(q) = \operatorname{gr}(f) - 1$: trovare una radice
permette di abbassare il grado del problema.
\end{osservazione}

\begin{definizione}
Sia $f \neq 0$. Una radice $\alpha$ ha \emph{molteplicità} (algebrica)
$m \geq 1$ se $(t - \alpha)^m \mid f$ ma $(t - \alpha)^{m+1} \nmid f$, cioè se
$f = (t - \alpha)^m\, g$ con $g(\alpha) \neq 0$. Le radici di molteplicità $1$
si dicono \emph{semplici}.
\end{definizione}

\subsection{Ricerca delle radici e scomposizione}

\begin{osservazione}
Se $f = a_n t^n + \dots + a_1 t + a_0$ ha coefficienti \emph{interi} e
$\frac{p}{q}$ (frazione ridotta) è una radice razionale di $f$, allora
$p \mid a_0$ e $q \mid a_n$. In particolare, se $f$ è monico, le radici
razionali sono intere e dividono $a_0$: basta provare i divisori di $a_0$
(con entrambi i segni).
\end{osservazione}

\begin{esempio}[radici semplici]
Scomponiamo $f(t) = t^3 - 2t^2 - 5t + 6$. Le possibili radici intere sono i
divisori di $6$: $\pm 1, \pm 2, \pm 3, \pm 6$. Si trova
\[
  f(1) = 1 - 2 - 5 + 6 = 0 ,
\]
quindi $1$ è radice e $(t - 1) \mid f$. Per trovare il quoziente si può
usare la regola di Ruffini: si riportano i coefficienti di $f$ e si procede
da sinistra a destra, moltiplicando per $\alpha = 1$ e sommando.
\[
\begin{array}{r|rrrr}
   & 1 & -2 & -5 & 6 \\
 1 &   &  1 & -1 & -6 \\ \hline
   & 1 & -1 & -6 & \;0
\end{array}
\]
L'ultimo numero è il resto ($=f(1)=0$), gli altri sono i coefficienti del
quoziente: $f = (t - 1)(t^2 - t - 6)$. Il fattore di secondo grado si
scompone cercando due numeri con somma $-1$ e prodotto $-6$:
$t^2 - t - 6 = (t - 3)(t + 2)$. Quindi
\[
  f(t) = (t - 1)(t - 3)(t + 2),
\]
e le radici sono $1$, $3$, $-2$, tutte semplici.
\end{esempio}

\begin{esempio}[radice multipla]
Sia $f(t) = t^3 - 3t^2 + 4$. I divisori di $4$ sono $\pm 1, \pm 2, \pm 4$ e
\[
  f(2) = 8 - 12 + 4 = 0 , \qquad f(-1) = -1 - 3 + 4 = 0 .
\]
Dividendo per $t - 2$ (Ruffini con $\alpha = 2$ sui coefficienti
$1, -3, 0, 4$) si ottiene $f = (t - 2)(t^2 - t - 2)$ e poi
$t^2 - t - 2 = (t - 2)(t + 1)$. Dunque
\[
  f(t) = (t - 2)^2 (t + 1) .
\]
La radice $2$ ha molteplicità $2$, la radice $-1$ ha molteplicità $1$.
\end{esempio}

\section{Ricorsione lineare}

\begin{definizione}
Una successione $(a_n)_{n \geq 0}$ è \emph{ricorsiva lineare omogenea di grado
$k$} (a coefficienti costanti) se esistono $c_1, \dots, c_k \in \mathbb{R}$,
con $c_k \neq 0$, tali che per ogni $n \geq k$
\[
  a_n = c_1\, a_{n-1} + c_2\, a_{n-2} + \dots + c_k\, a_{n-k} ,
\]
e se sono assegnati i primi $k$ valori $a_0, \dots, a_{k-1}$ (\emph{condizioni
iniziali}). Il polinomio
\[
  p(t) = t^k - c_1\, t^{k-1} - c_2\, t^{k-2} - \dots - c_{k-1}\, t - c_k
\]
si dice \emph{polinomio caratteristico} della ricorrenza.
\end{definizione}

``Lineare'' significa che $a_n$ è combinazione lineare dei termini precedenti,
``omogenea'' che non compaiono termini aggiuntivi che non dipendono dalla
successione.

\begin{osservazione}[perché il polinomio caratteristico]
Cerchiamo soluzioni della forma $a_n = \alpha^n$ con $\alpha \neq 0$.
Sostituendo nella ricorrenza e dividendo per $\alpha^{n-k}$ si ottiene
\[
  \alpha^k = c_1\, \alpha^{k-1} + \dots + c_k
  \iff p(\alpha) = 0 .
\]
Quindi $\alpha^n$ risolve la ricorrenza se e solo se $\alpha$ è una radice del
polinomio caratteristico. I teoremi seguenti dicono come combinare queste
soluzioni per soddisfare anche le condizioni iniziali.
\end{osservazione}

\begin{esempio}
\begin{itemize}
  \item $a_n = 3\, a_{n-1}$ ha grado $1$ e $p(t) = t - 3$; con $a_0$ dato si ha
        $a_n = a_0\, 3^n$.
  \item Fibonacci: $F_n = F_{n-1} + F_{n-2}$, con $F_0 = 0$, $F_1 = 1$. Grado
        $2$, $p(t) = t^2 - t - 1$.
  \item $a_n = 5\, a_{n-1} - 6\, a_{n-2}$ ha grado $2$ e
        $p(t) = t^2 - 5t + 6$.
  \item $a_n = 6 a_{n-1} - 11 a_{n-2} + 6 a_{n-3}$ ha grado $3$ e
        $p(t) = t^3 - 6t^2 + 11t - 6$.
\end{itemize}
\end{esempio}

\subsection{Teorema A: radici distinte}

\begin{teorema}[A]
Sia $(a_n)$ una successione ricorsiva lineare omogenea di grado $k$ il cui
polinomio caratteristico $p$ ha $k$ radici \emph{distinte}
$\alpha_1, \dots, \alpha_k$. Allora esistono costanti
$c_1, \dots, c_k$ tali che per ogni $n \geq 0$
\[
  a_n = c_1\, \alpha_1^{\,n} + c_2\, \alpha_2^{\,n} + \dots + c_k\, \alpha_k^{\,n} .
\]
Le costanti $c_i$ sono univocamente determinate dalle condizioni iniziali.
\end{teorema}

(Senza dimostrazione.)

\paragraph{Procedura per trovare la forma chiusa.}
\begin{enumerate}
  \item Scrivere il polinomio caratteristico $p(t)$.
  \item Trovare le sue $k$ radici distinte $\alpha_1, \dots, \alpha_k$
        (con il teorema del fattore).
  \item Scrivere $a_n = c_1 \alpha_1^{\,n} + \dots + c_k \alpha_k^{\,n}$.
  \item Imporre le condizioni iniziali $a_0, \dots, a_{k-1}$: si ottiene un
        sistema lineare di $k$ equazioni nelle incognite $c_1, \dots, c_k$.
  \item Risolvere il sistema e sostituire i valori trovati.
\end{enumerate}

\begin{esempio}
Sia $a_n = 5\, a_{n-1} - 6\, a_{n-2}$ con $a_0 = 0$, $a_1 = 1$. Il polinomio
caratteristico è $p(t) = t^2 - 5t + 6 = (t - 2)(t - 3)$, con radici distinte
$2$ e $3$. Quindi
\[
  a_n = c_1\, 2^n + c_2\, 3^n .
\]
Le condizioni iniziali danno il sistema
\[
  \begin{cases}
    a_0 = c_1 + c_2 = 0 \\
    a_1 = 2c_1 + 3c_2 = 1
  \end{cases}
  \quad\Longrightarrow\quad c_2 = 1,\ c_1 = -1 ,
\]
perciò
\[
  a_n = 3^n - 2^n .
\]
Verifica: $a_2 = 5 \cdot 1 - 6 \cdot 0 = 5 = 3^2 - 2^2$ e
$a_3 = 5 \cdot 5 - 6 \cdot 1 = 19 = 3^3 - 2^3$.
\end{esempio}

\begin{esempio}[Fibonacci]
Per $F_n = F_{n-1} + F_{n-2}$, $F_0 = 0$, $F_1 = 1$, il polinomio
caratteristico è $t^2 - t - 1$, con radici distinte
\[
  \varphi = \frac{1 + \sqrt{5}}{2} , \qquad \psi = \frac{1 - \sqrt{5}}{2} .
\]
Allora $F_n = c_1\, \varphi^n + c_2\, \psi^n$ e
\[
  \begin{cases}
    c_1 + c_2 = 0 \\
    c_1\, \varphi + c_2\, \psi = 1
  \end{cases}
  \quad\Longrightarrow\quad
  c_1 (\varphi - \psi) = 1 ,\quad \varphi - \psi = \sqrt{5}
  \quad\Longrightarrow\quad c_1 = \frac{1}{\sqrt{5}},\ c_2 = -\frac{1}{\sqrt{5}} .
\]
Si ottiene la formula di Binet:
\[
  F_n = \frac{1}{\sqrt{5}}\left[\left(\frac{1 + \sqrt{5}}{2}\right)^{n}
        - \left(\frac{1 - \sqrt{5}}{2}\right)^{n}\right] .
\]
\end{esempio}

\begin{esempio}[grado 3]
Sia $a_n = 6 a_{n-1} - 11 a_{n-2} + 6 a_{n-3}$ con $a_0 = 3$, $a_1 = 6$,
$a_2 = 14$. Il polinomio caratteristico è
$p(t) = t^3 - 6t^2 + 11t - 6$. Provando i divisori di $6$ si trova
$p(1) = 0$, e con Ruffini $p(t) = (t - 1)(t^2 - 5t + 6) = (t-1)(t-2)(t-3)$.
Le radici $1, 2, 3$ sono distinte, quindi
\[
  a_n = c_1 + c_2\, 2^n + c_3\, 3^n ,
  \qquad
  \begin{cases}
    c_1 + c_2 + c_3 = 3 \\
    c_1 + 2c_2 + 3c_3 = 6 \\
    c_1 + 4c_2 + 9c_3 = 14
  \end{cases}
\]
Sottraendo la prima equazione dalla seconda e la seconda dalla terza si ha
$c_2 + 2c_3 = 3$ e $2c_2 + 6c_3 = 8$, da cui $c_3 = 1$, $c_2 = 1$, $c_1 = 1$.
Dunque
\[
  a_n = 1 + 2^n + 3^n .
\]
\end{esempio}

\subsection{Teorema B: radici con molteplicità}

\begin{teorema}[B]
Sia $(a_n)$ una successione ricorsiva lineare omogenea di grado $k$ il cui
polinomio caratteristico si scompone come
\[
  p(t) = (t - \alpha_1)^{m_1} (t - \alpha_2)^{m_2} \cdots (t - \alpha_r)^{m_r},
  \qquad m_1 + m_2 + \dots + m_r = k ,
\]
con $\alpha_1, \dots, \alpha_r$ distinte e $m_i$ la molteplicità di
$\alpha_i$. Allora
\[
  a_n = \sum_{i=1}^{r}
        \bigl(c_{i,0} + c_{i,1}\, n + \dots + c_{i,m_i - 1}\, n^{m_i - 1}\bigr)\,
        \alpha_i^{\,n} ,
\]
con $k$ costanti $c_{i,j}$ univocamente determinate dalle condizioni iniziali.
\end{teorema}

(Senza dimostrazione.)

\begin{osservazione}
Ogni radice $\alpha$ di molteplicità $m$ contribuisce con $m$ termini,
\[
  \alpha^n,\ n\,\alpha^n,\ n^2\alpha^n,\ \dots,\ n^{m-1}\alpha^n ,
\]
in totale $k$ costanti, quante le condizioni iniziali. Se tutte le molteplicità
sono $1$ si ritrova il Teorema A.
\end{osservazione}

\begin{esempio}[radice doppia]
Sia $a_n = 4\, a_{n-1} - 4\, a_{n-2}$ con $a_0 = 1$, $a_1 = 0$. Il polinomio
caratteristico è $p(t) = t^2 - 4t + 4 = (t - 2)^2$: un'unica radice $2$ di
molteplicità $2$. Quindi
\[
  a_n = (c_0 + c_1\, n)\, 2^n .
\]
Dalle condizioni iniziali:
\[
  a_0 = c_0 = 1 , \qquad a_1 = 2\,(c_0 + c_1) = 0 \Rightarrow c_1 = -1 ,
\]
perciò
\[
  a_n = (1 - n)\, 2^n .
\]
Verifica: $a_2 = 4 \cdot 0 - 4 \cdot 1 = -4 = (1-2)\cdot 2^2$ e
$a_3 = 4\cdot(-4) - 4 \cdot 0 = -16 = (1-3)\cdot 2^3$.
\end{esempio}

\begin{esempio}[radici di molteplicità diversa]
Sia $a_n = 5\, a_{n-1} - 8\, a_{n-2} + 4\, a_{n-3}$ con $a_0 = 2$, $a_1 = 5$,
$a_2 = 13$. Il polinomio caratteristico è $p(t) = t^3 - 5t^2 + 8t - 4$.
Poiché $p(1) = 0$, con Ruffini si ottiene
\[
  p(t) = (t - 1)(t^2 - 4t + 4) = (t - 1)(t - 2)^2 ,
\]
cioè la radice $1$ ha molteplicità $1$ e la radice $2$ ha molteplicità $2$.
Per il Teorema B
\[
  a_n = c + (d + e\, n)\, 2^n ,
  \qquad
  \begin{cases}
    c + d = 2 \\
    c + 2d + 2e = 5 \\
    c + 4d + 8e = 13
  \end{cases}
\]
Sottraendo la prima equazione dalla seconda e la seconda dalla terza:
$d + 2e = 3$ e $2d + 6e = 8$, cioè $d + 3e = 4$. Quindi $e = 1$, $d = 1$,
$c = 1$ e
\[
  a_n = 1 + (1 + n)\, 2^n .
\]
Verifica: $a_3 = 5 \cdot 13 - 8 \cdot 5 + 4 \cdot 2 = 33 = 1 + 4 \cdot 8$.
\end{esempio}

\end{document}
